\documentclass[10pt,a4paper]{amsart}
\usepackage{amsmath,amssymb,amsthm,graphicx}
\usepackage{verbatim}
\newenvironment{reference}{\comment}{\endcomment}
%\newenvironment{reference}{}{}
\newenvironment{slogan}{\comment}{\endcomment}
\newenvironment{history}{\comment}{\endcomment}

% For commutative diagrams we use Xy-pic
\usepackage[all]{xy}

% We use 2cell for 2-commutative diagrams.
\xyoption{2cell}
% Theorem environments.
%
\theoremstyle{plain}
\newtheorem{theorem}[subsection]{Theorem}
\newtheorem{proposition}[subsection]{Proposition}
\newtheorem{lemma}[subsection]{Lemma}

\theoremstyle{definition}
\newtheorem{definition}[subsection]{Definition}
\newtheorem{example}[subsection]{Example}
\newtheorem{exercise}[subsection]{Exercise}
\newtheorem{situation}[subsection]{Situation}

\theoremstyle{remark}
\newtheorem{remark}[subsection]{Remark}
\newtheorem{remarks}[subsection]{Remarks}

\numberwithin{equation}{subsection}

% Macros
%
\def\lim{\mathop{\mathrm{lim}}\nolimits}
\def\colim{\mathop{\mathrm{colim}}\nolimits}
\def\Spec{\mathop{\mathrm{Spec}}}
\def\Hom{\mathop{\mathrm{Hom}}\nolimits}
\def\Ext{\mathop{\mathrm{Ext}}\nolimits}
\def\SheafHom{\mathop{\mathcal{H}\!\mathit{om}}\nolimits}
\def\SheafExt{\mathop{\mathcal{E}\!\mathit{xt}}\nolimits}
\def\Sch{\mathit{Sch}}
\def\Mor{\mathop{\mathrm{Mor}}\nolimits}
\def\Ob{\mathop{\mathrm{Ob}}\nolimits}
\def\Sh{\mathop{\mathit{Sh}}\nolimits}
\def\NL{\mathop{N\!L}\nolimits}
\def\CH{\mathop{\mathrm{CH}}\nolimits}
\def\proetale{{pro\text{-}\acute{e}tale}}
\def\etale{{\acute{e}tale}}
\def\QCoh{\mathit{QCoh}}
\def\Ker{\mathop{\mathrm{Ker}}}
\def\Im{\mathop{\mathrm{Im}}}
\def\Coker{\mathop{\mathrm{Coker}}}
\def\Coim{\mathop{\mathrm{Coim}}}

% Boxtimes
%
\DeclareMathSymbol{\boxtimes}{\mathbin}{AMSa}{"02}

%
% Macros for moduli stacks/spaces
%
\def\QCohstack{\mathcal{QC}\!\mathit{oh}}
\def\Cohstack{\mathcal{C}\!\mathit{oh}}
\def\Spacesstack{\mathcal{S}\!\mathit{paces}}
\def\Quotfunctor{\mathrm{Quot}}
\def\Hilbfunctor{\mathrm{Hilb}}
\def\Curvesstack{\mathcal{C}\!\mathit{urves}}
\def\Polarizedstack{\mathcal{P}\!\mathit{olarized}}
\def\Complexesstack{\mathcal{C}\!\mathit{omplexes}}
% \Pic is the operator that assigns to X its picard group, usage \Pic(X)
% \Picardstack_{X/B} denotes the Picard stack of X over B
% \Picardfunctor_{X/B} denotes the Picard functor of X over B
\def\Pic{\mathop{\mathrm{Pic}}\nolimits}
\def\Picardstack{\mathcal{P}\!\mathit{ic}}
\def\Picardfunctor{\mathrm{Pic}}
\def\Deformationcategory{\mathcal{D}\!\mathit{ef}}
\usepackage{cleveref}
\crefname{originallemma}{original lemma}{original lemmas}
\Crefname{originallemma}{Original lemma}{Original lemmas}
\makeatletter
\newlabel{chow-lemma-gysin-fundamental}{{0FBM}{source}{}{sourcepointer-1}{}}
\newlabel{chow-lemma-gysin-fundamental@cref}{{[originallemma][0][]0FBM}{[1][0][]source}{}{}{}}
\makeatother
\crefname{originallemma}{original lemma}{original lemmas}
\Crefname{originallemma}{Original lemma}{Original lemmas}
\makeatletter
\newlabel{chow-lemma-intersect-regularly-embedded}{{0FFE}{source}{}{sourcepointer-2}{}}
\newlabel{chow-lemma-intersect-regularly-embedded@cref}{{[originallemma][0][]0FFE}{[1][0][]source}{}{}{}}
\makeatother
\crefname{originalreference}{original reference}{original references}
\Crefname{originalreference}{Original reference}{Original references}
\makeatletter
\newlabel{lemma-class-diagonal}{{0FHH}{source}{}{sourcepointer-3}{}}
\newlabel{lemma-class-diagonal@cref}{{[originalreference][0][]0FHH}{[1][0][]source}{}{}{}}
\makeatother
\crefname{originalreference}{original reference}{original references}
\Crefname{originalreference}{Original reference}{Original references}
\makeatletter
\newlabel{lemma-degrees-cycles}{{0FHR}{source}{}{sourcepointer-4}{}}
\newlabel{lemma-degrees-cycles@cref}{{[originalreference][0][]0FHR}{[1][0][]source}{}{}{}}
\makeatother
\crefname{originalreference}{original reference}{original references}
\Crefname{originalreference}{Original reference}{Original references}
\makeatletter
\newlabel{lemma-push-unit}{{0FHG}{source}{}{sourcepointer-5}{}}
\newlabel{lemma-push-unit@cref}{{[originalreference][0][]0FHG}{[1][0][]source}{}{}{}}
\makeatother
\crefname{originalreference}{original reference}{original references}
\Crefname{originalreference}{Original reference}{Original references}
\makeatletter
\newlabel{lemma-pushforward}{{03M9}{source}{}{sourcepointer-6}{}}
\newlabel{lemma-pushforward@cref}{{[originalreference][0][]03M9}{[1][0][]source}{}{}{}}
\makeatother
\crefname{originallemma}{original lemma}{original lemmas}
\Crefname{originallemma}{Original lemma}{Original lemmas}
\makeatletter
\newlabel{morphisms-lemma-differentials-diagonal}{{08S2}{source}{}{sourcepointer-7}{}}
\newlabel{morphisms-lemma-differentials-diagonal@cref}{{[originallemma][0][]08S2}{[1][0][]source}{}{}{}}
\makeatother
\crefname{originalreference}{original reference}{original references}
\Crefname{originalreference}{Original reference}{Original references}
\makeatletter
\newlabel{remark-replace-cup-product}{{0FHD}{source}{}{sourcepointer-8}{}}
\newlabel{remark-replace-cup-product@cref}{{[originalreference][0][]0FHD}{[1][0][]source}{}{}{}}
\makeatother
\title{The self intersection of the diagonal, the Euler characteristic and the top Chern class}
\author{The Stacks Project Authors}
\date{stacks-project:weil.tex:local-snapshot-20261012; GNU FDL 1.2 or later}
\begin{document}
\maketitle
\noindent\textit{Editorial note.} All delivered mathematical text is the author's own, copied from the pinned Stacks Project chapter file weil.tex; no statement or proof step was rewritten, reordered or removed.\par\medskip
\noindent\textit{Editorial note.} The counted proof is the author's complete argument: the identification of the self intersection of the diagonal with the top Chern class of the tangent sheaf via the conormal sheaf, and the computation through the degree map and the Kunneth decomposition of the class of the diagonal.\par\medskip
\noindent\textit{Editorial note.} The label lemma-square-diagonal names the counted statement, and the proof environment immediately following it in the source is delivered verbatim as its proof.\par\medskip
\noindent\textit{Editorial note.} The counted statement is carried by the author's own lemma environment in the chapter source.\par\medskip
\noindent\textit{Editorial note.} Pre-open gate: python operations/agent-authoring/tools/precard\_gate.py --source weil.tex --label lemma-square-diagonal --start 3537 --end 3588 --local-id 28138 returned RESULT=OK (L28138 free, occupancy FREE, public interval clear). Pre-assembly tag self-check: canonical tag resolved and not present in reference/public-items.jsonl.\par\medskip
\noindent\textit{Editorial note.} The statement and proof are copied from the pinned chapter file; the cited inputs are kept as external pointers with their canonical Stacks tags.\par\medskip
\noindent\textit{Allowed standard prerequisites (editorial).}\begin{itemize}
\item The self intersection of the diagonal computed by the Gysin map. Chow Homology, the lemma expressing the intersection product of the diagonal with itself through the Gysin map and the lemma identifying the Gysin map with capping by the top Chern class of the normal sheaf. Produces the identification of the self intersection of the diagonal with the top Chern class of the tangent sheaf.
\item The conormal sheaf of the diagonal is the sheaf of differentials. Morphisms, the lemma identifying the conormal sheaf of the diagonal with the sheaf of differentials. Identifies the normal sheaf of the diagonal with the tangent sheaf.
\item Pairing with the class of the diagonal through the degree map. Same chapter, the lemma relating degrees of cycles to pairings of cycle classes, the lemmas on the pushforward of the unit and on pushforwards, and the lemma describing the class of the diagonal in the Kunneth decomposition. Computes the self intersection through the degree map and the Kunneth components of the class of the diagonal.
\end{itemize}\medskip
\ifdefined\hypertarget\hypertarget{sourcepointer-1}{}\fi
\noindent\textit{Source reference.} Original lemma 0FBM: sources/stacks-chow/tags.txt, source lines 19152--19152. Prerequisite pointer invoked at proof line(s) 3556; it names a standard fact of a sibling chapter that the counted argument uses, and every step of the counted proof built on it is retained verbatim in the delivered body.\par\medskip
\ifdefined\hypertarget\hypertarget{sourcepointer-2}{}\fi
\noindent\textit{Source reference.} Original lemma 0FFE: sources/stacks-chow/tags.txt, source lines 19284--19284. Prerequisite pointer invoked at proof line(s) 3551; it names a standard fact of a sibling chapter that the counted argument uses, and every step of the counted proof built on it is retained verbatim in the delivered body.\par\medskip
\ifdefined\hypertarget\hypertarget{sourcepointer-3}{}\fi
\noindent\textit{Source reference.} Original reference 0FHH: sources/stacks-weil/pinned-original-source.tex, source lines 2861--2861. Prerequisite pointer invoked at proof line(s) 3579; it names a standard fact of the same chapter that the counted argument uses, and every step of the counted proof built on it is retained verbatim in the delivered body.\par\medskip
\ifdefined\hypertarget\hypertarget{sourcepointer-4}{}\fi
\noindent\textit{Source reference.} Original reference 0FHR: sources/stacks-weil/pinned-original-source.tex, source lines 3505--3505. Prerequisite pointer invoked at proof line(s) 3564; it names a standard fact of the same chapter that the counted argument uses, and every step of the counted proof built on it is retained verbatim in the delivered body.\par\medskip
\ifdefined\hypertarget\hypertarget{sourcepointer-5}{}\fi
\noindent\textit{Source reference.} Original reference 0FHG: sources/stacks-weil/pinned-original-source.tex, source lines 2848--2848. Prerequisite pointer invoked at proof line(s) 3576; it names a standard fact of the same chapter that the counted argument uses, and every step of the counted proof built on it is retained verbatim in the delivered body.\par\medskip
\ifdefined\hypertarget\hypertarget{sourcepointer-6}{}\fi
\noindent\textit{Source reference.} Original reference 03M9: sources/stacks-weil/pinned-original-source.tex, source lines 2660--2660. Prerequisite pointer invoked at proof line(s) 3577; it names a standard fact of the same chapter that the counted argument uses, and every step of the counted proof built on it is retained verbatim in the delivered body.\par\medskip
\ifdefined\hypertarget\hypertarget{sourcepointer-7}{}\fi
\noindent\textit{Source reference.} Original lemma 08S2: sources/stacks-chow/tags.txt, source lines 11115--11115. Prerequisite pointer invoked at proof line(s) 3558; it names a standard fact of a sibling chapter that the counted argument uses, and every step of the counted proof built on it is retained verbatim in the delivered body.\par\medskip
\ifdefined\hypertarget\hypertarget{sourcepointer-8}{}\fi
\noindent\textit{Source reference.} Original reference 0FHD: sources/stacks-weil/pinned-original-source.tex, source lines 2742--2742. Prerequisite pointer invoked at proof line(s) 3581; it names a standard fact of the same chapter that the counted argument uses, and every step of the counted proof built on it is retained verbatim in the delivered body.\par\medskip
\begin{lemma}
\label{lemma-square-diagonal}
Assume given (D0), (D1), (D2), and (D3) satisfying (A), (B), and (C).
Let $X$ be a nonempty smooth projective scheme over $k$ which is
equidimensional of dimension $d$. We have
$$
\sum\nolimits_i (-1)^i\dim_F H^i(X) =
\deg(\Delta \cdot \Delta) = \deg(c_d(\mathcal{T}_{X/k}))
$$
\end{lemma}

\begin{proof}
Equality on the right. We have
$[\Delta] \cdot [\Delta] = \Delta_*(\Delta^![\Delta])$
(Chow Homology, Lemma \ref{chow-lemma-intersect-regularly-embedded}).
Since $\Delta_*$ preserves degrees of $0$-cycles it suffices to compute
the degree of $\Delta^![\Delta]$. The class $\Delta^![\Delta]$ is given
by capping $[\Delta]$ with
the top Chern class of the normal sheaf of $\Delta \subset X \times X$
(Chow Homology, Lemma \ref{chow-lemma-gysin-fundamental}).
Since the conormal sheaf of $\Delta$ is $\Omega_{X/k}$
(Morphisms, Lemma \ref{morphisms-lemma-differentials-diagonal})
we see that the normal sheaf is equal to the tangent sheaf
$\mathcal{T}_{X/k} = \SheafHom_{\mathcal{O}_X}(\Omega_{X/k}, \mathcal{O}_X)$
as desired.

\medskip\noindent
Equality on the left. By Lemma \ref{lemma-degrees-cycles} we have
\begin{align*}
\deg([\Delta] \cdot [\Delta])
& =
\int_{X \times X} \gamma([\Delta]) \cup \gamma([\Delta]) \\
& =
\int_{X \times X} \Delta_*1 \cup \gamma([\Delta]) \\
& =
\int_{X \times X} \Delta_*(\Delta^*\gamma([\Delta])) \\
& =
\int_X \Delta^*\gamma([\Delta])
\end{align*}
We have used Lemmas \ref{lemma-push-unit} and
\ref{lemma-pushforward}.
Write $\gamma([\Delta]) = \sum  e_{i, j} \otimes e'_{2d - i , j}$
as in Lemma \ref{lemma-class-diagonal}.
Recalling that $\Delta^*$ is given by cup product
(Remark \ref{remark-replace-cup-product}) we obtain
$$
\int_X \sum\nolimits_{i, j} e_{i, j} \cup e'_{2d - i, j} =
\sum\nolimits_{i, j} \int_X e_{i, j} \cup e'_{2d - i, j} =
\sum\nolimits_{i, j} (-1)^i = \sum (-1)^i\beta_i
$$
as desired.
\end{proof}

\end{document}
